\chapter{Applications of planar separator theorem}
\section{Matching algorithm}

\subsection{Algorithm}
In this section we will discuss an efficient algorithm for finding a maximum
matching in a planar graph. The maximum matching problem asks for a maximum
number of pairwise nonadjacent edges in a graph. Although there exists an
exact algorithm to find a maximum matching for planar graphs in $O(n \sqrt{n})$
time, we can based on theorems recalled in Chapter 2 exhibit an approximation
algorithm that runs in $O(n \log{n})$ time and whose worst-case ratio is 1 -
$O(1/ \sqrt{\log\log n})$. One can find it interesting why we cannot use Algorithm
$MISP$ directly and have worst-case ratio and running time guaranteed by Lemma~
\ref{lem:MISP-upperbound}. That is because the third property is not
satisfied. The counter example is graph $K_{1, n-1}$ whose matching size is 1
regardless of the size of the graph, and that property demands that the size of
the matching is a positive fraction of $n$. It is also really important to see
that even though $MISP$ algorithm was defined for problems with sets of vertices
it easily extendable to sets of edges which we can employ in this section.

\begin{algorithm}[H]
\caption{MATCHING procedure}
\label{alg:matching}
\begin{algorithmic}[1]
    \Procedure{MATCHING}{}
    \State \Call{REDUCE}{G, S(G)}
    \State \Return{S(G)}
    \EndProcedure
\end{algorithmic}
\end{algorithm}

\begin{algorithm}[H]
\caption{REDUCE procedure}
\label{alg:reduce}
\begin{algorithmic}[1]
    \Procedure{REDUCE}{$G, s(G)$}
        \State $N \gets |V(G)|$
        \State $v \gets \min_{u \in V(G)} \deg(u)$
        \If{$N \leq \log\log n$}
            \State Find exact matching $S^*(G)$ using any reasonable algorithm
            \State $S(G) \gets S^*(G)$
        \ElsIf{$\deg(v) < 3$}
            \If{$\deg(v) = 0$}
                \State $G' \gets G - v$
                \State \Call{REDUCE}{$G', S(G')$}
                \State $S(G) \gets S(G')$
            \ElsIf{$\deg(v) = 1$}
                \State $u \gets \text{vertex adjacent to } v$
                \State $G' \gets G - \{u, v\}$
                \State \Call{REDUCE}{$G', S(G')$}
                \State $S(G) \gets S(G') \cup \{ (v, u) \}$
            \ElsIf{$\deg(v) = 2$}
                \State $u, w \gets \text{neighbors of } v$
                \State $G' \gets \text{graph made from } G \text{ by identifying }
                v, u \text{ and } w$
                \State \Call{REDUCE}{$G', S(G')$}
                \If{$S(G') \text{ does not have an edge adjacent to } u \text{
                in } G$}
                    \State $S(G) \gets S(G') \cup \{(v, u)\}$
                \Else
                    \State $S(G) \gets S(G') \cup \{(v, w)\}$
                \EndIf
            \EndIf

        \Else
            \State Run Algorithm $MISP$ with $\epsilon =
            {\log\log n}/N$ on $G$
            \State $S(G) \gets \text{result of algorithm}$
        \EndIf
    \EndProcedure
\end{algorithmic}
\end{algorithm}

\begin{theorem}\label{thm:matching}
    The approximation algorithm MATCHING for the maximum matching problem on
    planar graphs has running time $O(n \log n)$ and worst-case ratio
    $1 - O(1/\sqrt{\log \log n}) $.
\end{theorem}

\begin{proof}
\subsection{Correctness}
We will start the proof to Theorem \ref{thm:matching} by proving that it returns
a matching with stated worst-case ratio. For that we will split the proof into
two parts. First we will show that reductions performed by REDUCE procedure
make the approximation error smaller than that of the graph before a reduction.
And secondly, we will show that in base cases the algorithm returns the correct
worst-case ratio. We will define the maximum matching of $G$ as $S^*(G)$ and the
matching produced by algorithm as $S(G)$. We will show that
\begin{align*}
    \frac{|S(G)|}{|S^*(G)|} \geq 1 - O\left( \frac{1}{\sqrt{\log\log n}} \right)
\end{align*}

\subsubsection{Reductions}
\begin{enumerate}
    \item $G$ has a vertex of degree 0. Let $G' = G - v$. Then we have that
        $S(G) = S(G')$, as there is no edge adjacent to $v$. For the same reason
        we have $S^*(G) = S^*(G')$, so the reduction does not affect the ratio.
    \item $G$ has a vertex of degree 1. Let $u$ be the vertex adjacent to $v$ in
        $G$ and $G' = G - \{u, v\}$. If $S(G')$ is a matching produced by the
        algorithm then clearly $|S(G)| = |S(G')| + 1$ as we can add $(u,v)$ to
        the matching. Also since $S^*(G)$ has exactly one edge adjacent to $u$
        then by removing it we can reduce the matching by at most one, so we
        have $S^*(G) \leq S^*(G') + 1$. Now we can derive:
        $$
         \frac{|S(G)|}{|S^*(G)|} \geq
         \frac{|S(G')| + 1}{|S^*(G')| + 1} \geq
         \frac{|S(G')|}{|S^*(G')|}
        $$
        where the last equality stems from properties of rational numbers.
    \item $G$ has a vertex of degree 2. Let $u, w$ be vertices adjacent to $v$
        and $G'$ be a graph obtained from $G$ by identifying those three
        vertices into single vertex $v'$. Let $S(G) = S(G') \cup \{(u,v)\}$ if no
        edge adjacent to $u$ in $G$ was selected to $S(G')$. Otherwise let
        $S(G) = S(G') \cup (w,v)$. Clearly we have $|S(G)| = |S(G')| + 1$. Now
        if the matching $S^*(G)$ contains either $(v,u)$ or $(v,w)$, let's say
        $(v,u)$ then by removing it by contraction to $v'$, $S^*(G')$ is a
        correct matching of size at most one less than the latter. In the other
        case if both $u$ and $w$ were incident to edges in $S^*(G)$, then the
        matching $S^*(G')$ has to select one of them, so the property is
        conserved. In both cases we have $|S^*(G)| \leq |S^*(G')| + 1$. So again

        $$
         \frac{|S(G)|}{|S^*(G)|} \geq
         \frac{|S(G')| + 1}{|S^*(G')| + 1} \geq
         \frac{|S(G')|}{|S^*(G')|}
        $$
\end{enumerate}
We can see that all of the reductions create a graph if a worst-case ratio no
worse than the original graph $G$.
\subsubsection{Base cases}
We have two base cases to discuss here:
\begin{enumerate}
    \item If $N \leq \log{\log{n}}$, so we reduced our graph to a small fraction
        of the original one, we apply an exact matching algorithm, so
        $\frac{|S(G)|}{|S^*(G)|} = 1$ (if the denominator is 0, then we set this
        ratio to 1)
    \item If $N \geq \log\log n$ and $\forall_{v \in V(G)} \, \deg(v) \geq 3$
        we apply $MISP$ algorithm. To prove its correctness we have to check conditions
        from Lemma~\ref{lem:MISP-upperbound}:
        \begin{enumerate}
            \item Satisfied, as $S(G)$ is a matching so a subgraph induced by it
                satisfies the property that no two edges can have a common
                vertex
            \item Let $S^*(G)$ be the maximum matching. By removing vertices of
                a separator $C$ at most $|C|$ edges from the matching have
                incident vertices in $C$. So in the induced graph $G - C$ the
                maximum matching is of size at least $S^*(G) - |C|$. And this is
                also the size of $S(G)$ as we run the exact algorithm for
                matching here. So $|S(G)| \geq S^*(G) - |C|$, by rearranging we
                have $|S^*(G)| - |S(G)| \leq |C| = O(|C|)$, so the error is bound
                by size of $C$.
            \item Here we will use the theorem from \cite{nishizeki1979lower}
                which guarantees that if $G$ is planar and has minimum degree of
                at least 3, then $|S^*(G)| \geq N/3$, so the optimal solution is a
                positive fraction of $N$.
            \item We can easily check in linear time if the set found by the
                algorithm is the matching
        \end{enumerate}
        We can see that all the conditions are satisfied, so by setting
        $\epsilon = \log{\log{n}}/N$ we have the desired worst-case error by
        Lemma~\ref{lem:MISP-upperbound}.
\end{enumerate}

\subsection{Time complexity}
At first we should discuss the time complexity of reductions. By keeping a graph
as adjacency lists, so that each vertex not only has a previous reference and
next reference but also the cross reference to other linked lists, we can
perform a deletion in $O(\deg(v))$. And as each vertex is deleted only once, we
can do that in amortized time $O(1)$. Another operation is identification of
vertices which can be performed in total time $O(n \log{n})$
\cite{chibaApplicationsLiptonTarjans1981}. To do that we keep an additional
adjacency matrix and rewiring edges from the vertex of a smaller degree to one
of higher. As each vertex is present in $O(\log{n})$ such rewirings, we get the
needed time complexity.
Now we can look at running time of base cases. If $N \leq \log\log n$ we use
an exact algorithm to find the matching. If its time complexity is at most
$O(n2^N)$, then it is enough for our desired time complexity. And if $N \geq
\log\log n$, then from Lemma~\ref{lem:MISP-upperbound} we have guaranteed
$O(\max \{ N \log{N}, N 2^{ \epsilon N} \}) = O(N \log{n}) = O(n \log{n})$.

\end{proof}

\section{Vertex cover algorithm}
Another problem we will look at is Vertex Cover which is NP-complete. In
\cite{chibaApplicationsLiptonTarjans1981} we can find an algorithm very similar
to the one for matching problem discussed in a previous section. However for
VC another interesting approach may be found in
\cite{bar-yehudaApproximatingVertexCover1982} which will be shown in this
section.
At first we shall define what a vertex cover is.
\begin{defn}\label{def:vertex-cover}
    Let $G = (V, E)$ be a graph. We call $VC \subseteq V$ a vertex cover if for each 
    $e \in E$ one of its endpoints is in $VC$.
\end{defn}
Vertex Cover Problem asks for a minimum $VC$ in $G$. Another related problem is the
Independent Set Problem that asks for a maximum independent set.
\begin{defn}
    Let $G = (V, E)$ be a graph. We call $I \subseteq V$ and independent set if for 
    every two $v, u \in I$, $(u, v) \notin E$.
\end{defn}
It is clear that those problems are complementary, so solving Independent Set Problem
gives a solution to Vertex Cover Problem as $VC \cup I = V$. In \cite{liptonApplicationsPlanarSeparator1977}
authors give an approximation algorithm (which is essentially MISP algorithm discussed in Chapter 2) to find
a maximum independent set in $O(n \log n)$ time with
relative error $O(1/\sqrt{\log \log n})$, which means that for some constant $c$:
$$
\frac{|I^*| - |I|}{|I^*|} \leq \frac{c}{\sqrt{\log \log n}}
$$
But we know that $|VC^*| + |I^*| = n$, so after some simple arithmetics:
$$
\frac{|VC| - |VC^*|}{|VC^*|} \leq \frac{c (n/|VC^*| - 1)}{\sqrt{\log \log n}}
$$
This means that in order to use approximation scheme from \cite{liptonApplicationsPlanarSeparator1977}
for independent set to solve VC Problem we will have to show that $n/|VC^*|$ converges
to some $\epsilon_0 > 1$.
\\
To do that we will have to pre-process graph $G$ and apply approximation scheme afterwards.
At start we shall start with a few definitions introduced by authors of
\cite{bar-yehudaApproximatingVertexCover1982}.


\begin{defn}\label{def:rosary}
Let $G = (V, E)$ be a simple planar graph. A rosary in $G$ is a simple circuit
of even length $2l$, denoted by $v_1 - v_2 - \dots - v_{2l} - v_1$, such that the
degree of every even-indexed vertex is exactly two
for all $i = 1, 2, \dots, l$.
\end{defn}
\begin{defn}\label{def:bead}
Let $G = (V, E)$ be a simple planar graph. A vertex $v \in V$ of degree two is
called a bead if it plays the role of an even-indexed vertex $v_{2i}$ in
some rosary in $G$.
\end{defn}
\begin{defn}
    Let $G = (V, E)$ be a simple planar graph. We will call it proper if the degree
    of all vertices is at least 2 and the set $Y_B$ of all its beads satisfies
$|Y_B| < 1/6 |V|$.
\end{defn}
\noindent
In \cite{bar-yehudaApproximatingVertexCover1982} authors prove a very useful
theorem regarding proper graphs.
\begin{theorem}\label{thm:proper}
    Let $G = (V, E)$ be a simple planar graph with all vertices having degree at
    least 2. Then we have:
    $$
    |VC^*| > \frac{1}{5}(|V| - |Y_B|)
    $$
\end{theorem}
\noindent
If $G$ is proper then by Theorem~\ref{thm:proper} we have $|V|/|VC^*| < 6$. And that is exactly 
what we need to apply the approximation scheme. So the preparatory algorithm from
\cite{bar-yehudaApproximatingVertexCover1982} will reduce $G$ to a proper graph. The general
idea of that algorithm is to at first get rid of all vertices with degree less than 2 and later 
reduce the size of set of beads. We shall see exact steps in listing for Algorithm~\ref{alg:prep}.

\begin{algorithm}[H]
\caption{Preparatory algorithm}
\label{alg:prep}
\small
\begin{algorithmic}[1]
    \Procedure{PREP}{$G$}
    \State $ VC \gets \emptyset$
    \State $ U \gets \emptyset $
    \State $ STOP \gets 0$
    \While{$STOP \neq 1$}
    \While{$\exists_{v \in V(\wtG(V \setminus U))} \deg(v) < 2$}
    \If{$v \in V \setminus U$, $ \deg(v) = 0$ in $\wtG(V \setminus U))$}
        \State Add $v$ to $U$
    \EndIf
    \If{$v \in V \setminus U$, $ \deg(v) = 1$ in $\wtG(V \setminus U))$ }
        \State Let $(v, u) \in E(\wtG(V \setminus U))$ be the edge of $v$
        \State Add $v, u$ to $U$ and $u$ to $VC$
    \EndIf
    \EndWhile

    \While{There are simple isolated circuits in $\wtG(V \setminus U)$}
    \State Let $v_1, \dots, v_k, v_1$ be a simple isolated circuit in $\wtG(V \setminus U)$
    \State $\forall$ odd $i \in \{1, \dots, k\}$ $\rightarrow$ add $v_i$ to $VC$
    \State $\forall$ $i \in \{1, \dots, k\}$ $\rightarrow$ add $v_i$ to $U$
    \EndWhile
    \State $B(X, Y, E') \gets \text{BIPARTITE($\wtG(V \setminus U)$)}$
    \If{$|Y| \geq 1/6 |V \setminus U|$}
        \State $VC \gets VC \cup X$
        \State $U \gets U \cup X \cup Y$
    \EndIf
    \If{$|Y| < 1/6 |V \setminus U|$ or $V \setminus U = \emptyset$}
        \State $STOP \gets 1$
    \EndIf
    \EndWhile
    \EndProcedure
\end{algorithmic}
\end{algorithm}

\subsubsection{Correctness of the algorithm}
We shall at first prove that Algorithm~\ref{alg:prep} outputs a minimum vertex cover
for $\wtG(U)$ and $\wtG(V \setminus U)$ which shall be a proper graph.
The algorithm has three phase that can be extracted:
\begin{enumerate}
    \item The first phase is deleting all the vertices with vertex degree less than 2.
        We can see that the that the covering of deleted induced subgraphs is optimal.
        \begin{itemize}
            \item If $\deg(v) = 0$, then it does not have an incident edge to it, so 
                the algorithms does not take it into the vertex cover which is optimal
                for the minimum vertex cover.
            \item If $\deg(v) = 1$, then it has to be picked to a minimum vertex cover or 
                its adjacent vertex $u$ to cover the incident edge $e$. The algorithm here 
                makes an optimal choice one again of picking $u$ to the vertex cover, as 
                it may have more incidents than only $e$ as opposed to $v$.
        \end{itemize}

    \item The second phase is deleting all of the isolated cycles. It is easy to see that 
        here the algorithm makes an optimal decision, as it adds to the vertex cover all the
        vertices with odd indices. It is the optimal decision, as the cycle has exactly $k$ edges
        to cover, so the minimum number of vertices to pick to cover the cycle 
        $\left \lceil{k/2}\right \rceil $, which is achieved by picking all the odd indices.
\begin{algorithm}[H]
\caption{BIPARTITE procedure}
\label{alg:bipartite}
\small
\begin{algorithmic}[1]
    \Procedure{BIPARTITE}{$H$ with no isolated circuits and $\deg(v) \geq 2$}
    \State Label all $v : \deg(v) > 2$ with "+"
    \State Add those vertices to queue $Q$
    \While{$Q$ not empty}
    \State $v \gets Q.popFront()$
    \State Let $L$ be all unlabelled vertices adjacent to $v$ in $H$
    \State $\forall v_L \in L$ $\rightarrow$ label opposite to $v$ (using "+" or "-")
    \State $\forall v_L \in L$ $\rightarrow$ $Q.add(v_L)$
    \EndWhile
    \State $X \gets \{ v : v \in V(H) \, \text{and labelled "+"}\}$
    \State $Y \gets \{ v : v \in V(H) \, \text{and labelled "-"}\}$
    \State $E' \gets \{(u, v) : u \in X, v \in Y, (u, v) \in E(H)\}$
    \State Let $B = (X, Y, E')$ be a bipartite graph
    \While{$\exists v \in B: \deg_B(v) < 2$}
        \State Remove $v$ and incident edges
    \EndWhile
    \State \Return $B$
    \EndProcedure
\end{algorithmic}
\end{algorithm}

    \item Now we shall focus on the third part of the algorithm, which is the construction of 
        a bipartite graph from the residual graph $\wtG(V \setminus U)$. The algorithm picks the whole
        vertex set $X$ to $VC$. To prove that it is optimal we have to show:
        \begin{lemma} \label{lem:vc-bipartie}
            Let $B(X, Y, E)$ be a bipartite graph in which for every $x \in X$, $\deg(x) \geq 2$ and 
            for every $y \in Y$, $\deg(y) \leq 2$.
        \end{lemma}
        \begin{proof}
            The proof follows by induction on $|Y|$. For the base case $|Y| = 0$, we have that $|X| = 0$
            (because the degree for each vertex of $x$ has to be at least 2) and $E = \emptyset$ which 
            trivially holds.
            \\ 
            Now for induction step suppose that the theorem holds for all $|Y| < m$, and let $B(X, Y, E)$
            be such a bipartite graph that $|Y| = m$. Let $y$ the vertex in $Y$ with minimum degree.
            If $\deg(y) = 0$ then delete $y$ from $B$ and the lemma is true by induction on the reduced graph, 
            so it is true for the original graph as well. If $\deg(y) = 1$, let $x$ by the adjacent vertex from $X$.
            Remove $x$, $y$ and all their incident edges. The set $X \setminus \{x\}$ is therefore the 
            minimum vertex cover of the bipartite graph $B(\{X \setminus x\}, \{Y \setminus y\}, E')$ by induction.
            So for the original graph add $x$ to the vertex cover, so that it covers all the removed incident edges.
            By doing that we have that $X$ is the minimum vertex cover of the original graph. If $\deg(y) = 2$, then 
            all vertices in $B(X, Y, E)$ have degrees at least 2, so there are cycles. Pick any simple cycle $C$. All the 
            vertices from $C$ that are in $X$ form the minimum vertex cover for that cycle. Remove $C$ and all the incident 
            edges. The minimum vertex cover for $B(X', Y', E')$ is $X'$ by induction, so with addition of the deleted 
            vertices from $C$ that formed the minimum vertex cover for $C$ we have that $X$ is the minimum vertex cover for
            the original graph $B(X, Y, E)$.
        \end{proof}
        Now we can see that by construction from the Algorithm~\ref{alg:bipartite}, all of the vertices with degree 
        greater than 2 land in set $X$ of the bipartite graph, so by Lemma~\ref{lem:vc-bipartie} $X$ is the minimum 
        vertex cover of $B(X, Y, E)$.
\end{enumerate}
Now we can observe that if Algorithm~\ref{alg:prep} halts with $V \setmius U = \emptyset$, then by the arguments that were 
just established, $VC$ is the minimum vertex cover of the whole graph, and the residual graph is empty. On the other hand
if the algorithm halts with the condition $|Y| < 1/6 |V \setminus U|$, then set $VC$ is the minimum vertex cover of $\wtG(U)$.
Moreover the degrees of all of the vertices is at least 2. The only thing that is left to prove is that the set of beads 
$Y_B$ satisfies $|Y_B| < 1/6 |V \setminus U|$.
\begin{lemma}
    If $y$ is a bead, then $y \in Y$.
\end{lemma}
\begin{proof}
    Let $(y_1, x_1, \dots, y_l, x_l)$ be a rosary in $H$ (input graph to Algorithm~\ref{alg:bipartite}). Since $H$
    contains no isolated circuits as a result of earlier steps for Algorithm~\ref{alg:prep}, some of the $x$ vertices 
    are of degree greater that 2 and are labelled $+$ in Algorithm~\ref{alg:bipartite}. All the other vertices in a rosary 
    are of degree 2 (so their neighbors are only other vertices from the rosary) and therefore their labelling cannot be
    affected by a vertex not from the rosary. It follows that every bead will finally get labelled $-$, so all of the beads will
    land in $Y$.
\end{proof}
By this lemma we have that $Y_B \subseteq Y$, so $\wtG(V \setminus U)$ is a proper graph.

\subsubsection{Time complexity}
We can show that the time complexity of the preparatory Algorithm~\ref{alg:prep} is $O(n)$. Construction of the 
bipartite graph, so the running time of Algorithm~\ref{alg:bipartite} is $O(n)$ as this is just a BFS. Moreover 
time spent on all of the transformations like deletion of vertices and isolated circuits is also $O(n)$. Note that
with every step of the outer while loop of Algorithm~\ref{alg:prep}, the number of vertices in $V \setminus U$ is 
reduced by a factor of $5/6$. Let $n_i = |V \setminus U|$ at i-th step of the outer while loop. Let $cn_i$ be the bound 
of the complexity of the i-th step. Since $n_i \leq (5/6)^i n$ then the total running time of Algorithm~\ref{alg:prep} is 
$$
\sum_{i=0}^{\infty} cn_i \leq \sum_{i=0}^{\infty} (5/6)^i n = 6 \cdot n \cdot c = O(n)
$$





